\documentclass[11pt]{article}
\usepackage[utf8]{inputenc}
\usepackage[T1]{fontenc}
\usepackage{amsmath,amssymb,amsthm}
\usepackage{graphicx}
\usepackage{hyperref}
\usepackage[export]{adjustbox}

\title{Algebra}
\author{Riccardo Di Michele}
\date{}

\begin{document}

\maketitle

\section{Insiemi}
Iniziamo con delle nozioni molto basilari. \\
Definiamo un \emph{\textbf{insieme}} come una raccolta di elementi. Le persone all'interno
di uno stadio sono un \emph{insieme}, gli snack all'interno della mia dispensa sono un
\emph{insieme} e, come vedremo, anche i numeri fanno parte di un \emph{insieme}. \\
Indicheremo quindi un insieme come: 

\[
  A = \{ 1, 2, 3, 4, 5 \}
\]

Il numero degli elementi all'interno di un insieme prende il nome di \textbf{cardinalità}
e può essere \emph{finito} o \emph{infinito}. Lo indicheremo come \(|A|\)

\end{document}